\documentclass[11pt,a4paper]{article}
\usepackage[margin=1in]{geometry}
\usepackage{amsmath,amssymb,amsthm}
\usepackage{algorithm}
\usepackage{algpseudocode}
\usepackage{booktabs}
\usepackage{hyperref}

\theoremstyle{definition}
\newtheorem{definition}{Definition}[section]
\newtheorem{theorem}{Theorem}[section]
\newtheorem{lemma}[theorem]{Lemma}
\newtheorem{remark}{Remark}[section]

\title{\textbf{Theoretical Foundations, Manifold Preservation, and Inverted Residual Mechanics in MobileNetV2}}
\author{Team Scaibu \\ \small Email: \texttt{jaydeep.9664920749@gmail.com} \\ \small Architecture and Core Engine Team}
\date{October 2026}

\begin{document}
\maketitle

\begin{abstract}
This document presents the formal mathematical specification, manifold preservation theory, and structural mechanics of Inverted Residuals and Linear Bottlenecks (MobileNetV2). Standard residual blocks compress channels at the bottleneck and place identity shortcuts across wide high-dimensional channels. Inverted residual blocks invert this design: shortcuts connect narrow low-dimensional bottlenecks, while intermediate spatial processing operates in an expanded high-dimensional space ($t \times C_{\text{in}}$). Crucially, the final projection omits non-linear activations to prevent manifold collapse in low-dimensional subspaces. We formulate the 3-stage transformation, prove linear manifold preservation, analyze computational complexities, trace a numerical example, and document memory-saving execution patterns.
\end{abstract}

\section{Notation and Mathematical Preliminaries}

Let $X \in \mathbb{R}^{C_{\text{in}} \times H_{\text{in}} \times W_{\text{in}}}$ denote the low-dimensional input activation tensor entering the inverted bottleneck block.

\begin{itemize}
    \item $C_{\text{in}}, C_{\text{out}} \in \mathbb{N}_{\ge 1}$: Input and output channel dimensions (bottleneck widths).
    \item $t \in \mathbb{R}_{\ge 1}$: Expansion multiplier factor ($t = 6$ typical).
    \item $C_{\text{exp}} = t \cdot C_{\text{in}}$: Expanded intermediate channel dimensionality.
    \item $W_{\text{exp}} \in \mathbb{R}^{C_{\text{exp}} \times C_{\text{in}}}$: $1 \times 1$ pointwise expansion projection weights.
    \item $W_{\text{dw}} \in \mathbb{R}^{C_{\text{exp}} \times K_h \times K_w}$: $3 \times 3$ depthwise spatial convolution filter weights.
    \item $W_{\text{proj}} \in \mathbb{R}^{C_{\text{out}} \times C_{\text{exp}}}$: $1 \times 1$ linear pointwise reduction projection weights.
    \item $\sigma_6(z) = \min(\max(0, z), 6)$: Bounded ReLU6 non-linearity.
    \item $Y \in \mathbb{R}^{C_{\text{out}} \times H_{\text{out}} \times W_{\text{out}}}$: Output activation tensor.
\end{itemize}

\begin{table}[h]
\centering
\caption{Summary of Mathematical Symbols for Inverted Bottlenecks}
\begin{tabular}{lll}
\toprule
\textbf{Symbol} & \textbf{Type / Domain} & \textbf{Description} \\
\midrule
$X$ & $\mathbb{R}^{C_{\text{in}} \times H_{\text{in}} \times W_{\text{in}}}$ & Input low-dimensional tensor \\
$t$ & $\mathbb{R}_{\ge 1}$ & Channel expansion ratio ($t = 6$) \\
$W_{\text{exp}}$ & $\mathbb{R}^{(t C_{\text{in}}) \times C_{\text{in}}}$ & Expansion $1 \times 1$ weights \\
$W_{\text{dw}}$ & $\mathbb{R}^{(t C_{\text{in}}) \times 3 \times 3}$ & Depthwise spatial filtering weights \\
$W_{\text{proj}}$ & $\mathbb{R}^{C_{\text{out}} \times (t C_{\text{in}})}$ & Linear projection $1 \times 1$ weights \\
$\sigma_6$ & $\mathbb{R} \to [0, 6]$ & Bounded ReLU6 activation operator \\
$Y$ & $\mathbb{R}^{C_{\text{out}} \times H_{\text{out}} \times W_{\text{out}}}$ & Output feature tensor \\
\bottomrule
\end{tabular}
\end{table}

\section{Problem Definition and Core Concepts}

\subsection{Intuitive Meaning}
Applying non-linear activations (such as ReLU) directly to low-dimensional representations irrevocably zeroes out negative coordinate subspaces, destroying informative manifold geometry. Inverted residuals expand channels prior to non-linear depthwise filtering, enabling rich non-linear feature extraction in high dimensions ($C_{\text{exp}}$) while maintaining a purely linear projection back into the narrow bottleneck ($C_{\text{out}}$) without non-linearity.

\subsection{Formal Mathematical Recurrences}

\begin{enumerate}
    \item \textbf{1x1 Expansion}:
    \begin{equation}
    Z_{\text{exp}} = \sigma_6(W_{\text{exp}} *_{1\times 1} X) \quad \in \mathbb{R}^{C_{\text{exp}} \times H_{\text{in}} \times W_{\text{in}}}
    \end{equation}
    \item \textbf{Depthwise Spatial Convolution}:
    \begin{equation}
    Z_{\text{dw}} = \sigma_6(W_{\text{dw}} *_{\text{dw}} Z_{\text{exp}}) \quad \in \mathbb{R}^{C_{\text{exp}} \times H_{\text{out}} \times W_{\text{out}}}
    \end{equation}
    \item \textbf{Linear Pointwise Projection}:
    \begin{equation}
    Z_{\text{proj}} = W_{\text{proj}} *_{1\times 1} Z_{\text{dw}} \quad \in \mathbb{R}^{C_{\text{out}} \times H_{\text{out}} \times W_{\text{out}}}
    \end{equation}
    \item \textbf{Residual Addition}:
    \begin{equation}
    Y = \begin{cases} X + Z_{\text{proj}} & \text{if } s = 1 \land C_{\text{in}} = C_{\text{out}} \\ Z_{\text{proj}} & \text{otherwise} \end{cases}
    \end{equation}
\end{enumerate}

\section{Algorithm Specification}

\begin{algorithm}[H]
\caption{Inverted Residual and Linear Bottleneck Forward Pass}
\label{alg:inverted_residual}
\begin{algorithmic}[1]
\Require Input $X \in \mathbb{R}^{C_{\text{in}} \times H_{\text{in}} \times W_{\text{in}}}$, weights $W_{\text{exp}}, W_{\text{dw}}, W_{\text{proj}}$, stride $s$, padding $p$.
\Ensure Output $Y \in \mathbb{R}^{C_{\text{out}} \times H_{\text{out}} \times W_{\text{out}}}$, intermediate $Z_{\text{exp}}, Z_{\text{dw}}$, residual flag.
\State $Z_{\text{exp}} \gets \sigma_6(\text{Conv1x1}(X, W_{\text{exp}}))$
\State $Z_{\text{dw}} \gets \sigma_6(\text{DepthwiseConv2D}(Z_{\text{exp}}, W_{\text{dw}}, s, p))$
\State $Z_{\text{proj}} \gets \text{Conv1x1}(Z_{\text{dw}}, W_{\text{proj}})$ \Comment{Linear projection with NO non-linearity}
\If{$s = 1 \land C_{\text{in}} = C_{\text{out}}$}
    \State $Y \gets X + Z_{\text{proj}}$, $\text{has\_res} \gets \text{true}$
\Else
    \State $Y \gets Z_{\text{proj}}$, $\text{has\_res} \gets \text{false}$
\EndIf
\State \Return $\{Y, Z_{\text{exp}}, Z_{\text{dw}}, \text{has\_res}\}$
\end{algorithmic}
\end{algorithm}

\section{Theoretical Guarantees and Manifold Preservation}

\begin{theorem}[Manifold Collapse Prevention by Linear Bottlenecks]
Let $\mathcal{M} \subset \mathbb{R}^d$ be a $d$-dimensional manifold embedded in $\mathbb{R}^d$. If a non-linear activation $\text{ReLU}(W x)$ is applied with $W \in \mathbb{R}^{d \times d}$, any points mapped to $W x \le 0$ collapse into low-dimensional boundaries of dimension $< d$. Conversely, a linear transformation $W_{\text{proj}} x$ with $\text{rank}(W_{\text{proj}}) = d$ preserves manifold dimensionality without topological collapse:
\begin{equation}
\dim(W_{\text{proj}}(\mathcal{M})) = \dim(\mathcal{M}) = d
\end{equation}
\end{theorem}

\begin{proof}
By the rank-nullity theorem, a linear map with full rank $d$ is an injective homeomorphism on its range, ensuring that zero volume collapse occurs. Bounded ReLU6 in the expanded dimension $C_{\text{exp}} = 6d$ performs rich feature transformations where the probability of all coordinates zeroing simultaneously is $(1/2)^{6d} \approx 0$, preserving complete manifold invertibility.
\end{proof}

\subsection{Limitations}
\begin{itemize}
    \item \textbf{Residual Condition}: The identity shortcut is strictly bypassed when stride $s > 1$ or when input channels do not match output channels.
\end{itemize}

\section{Complexity Analysis}

\subsection{Time Complexity}
\begin{itemize}
    \item \textbf{FLOP Total}: $H_{\text{in}} W_{\text{in}} C_{\text{in}} C_{\text{exp}} + H_{\text{out}} W_{\text{out}} C_{\text{exp}} K^2 + H_{\text{out}} W_{\text{out}} C_{\text{exp}} C_{\text{out}}$.
    \item \textbf{Total Time Complexity}: $\Theta(H_{\text{out}} \cdot W_{\text{out}} \cdot t C_{\text{in}} (C_{\text{in}} + K^2 + C_{\text{out}}))$ FLOPs.
\end{itemize}

\subsection{Space Complexity}
\begin{itemize}
    \item \textbf{Storage}: $\mathcal{O}(H_{\text{out}} W_{\text{out}} (C_{\text{exp}} + C_{\text{out}}))$ auxiliary memory.
\end{itemize}

\section{Exactness and Error Analysis}

The operation executes exact floating-point matrix multiplications and piecewise linear ReLU6 clampings.

\section{Worked Numerical Example}

Let $C_{\text{in}} = 1, C_{\text{exp}} = 2, C_{\text{out}} = 1, H = 1, W = 1, s = 1, p = 0$.
Input $X = [[[2.0]]]$.
Weights: $W_{\text{exp}} = [[1.0], [2.0]], W_{\text{dw}} = [[[1.0]], [[0.5]]], W_{\text{proj}} = [[0.5, 0.5]]$.

\begin{enumerate}
    \item \textbf{Expansion + ReLU6}:
    $Z_{\text{exp}} = [\sigma_6(1(2)), \sigma_6(2(2))] = [\sigma_6(2), \sigma_6(4)] = [2.0, 4.0]^T$.
    \item \textbf{Depthwise + ReLU6}:
    $Z_{\text{dw}} = [\sigma_6(1(2)), \sigma_6(0.5(4))] = [2.0, 2.0]^T$.
    \item \textbf{Linear Projection}:
    $Z_{\text{proj}} = 0.5(2.0) + 0.5(2.0) = 1.0 + 1.0 = 2.0$.
    \item \textbf{Residual Addition}:
    $Y = X + Z_{\text{proj}} = 2.0 + 2.0 = [[[4.0]]]$.
\end{enumerate}

\section{Numerical Considerations and Edge Cases}

\begin{itemize}
    \item \textbf{ReLU6 Saturation}: Bounding activations to $\le 6.0$ ensures optimal FP16 dynamic range preservation during mobile INT8/FP16 quantization.
\end{itemize}

\begin{thebibliography}{9}
\bibitem{sandler2018mobilenetv2} M.~Sandler, A.~Howard, M.~Zhu, A.~Zhmoginov, and L.~C.~Chen, ``MobileNetV2: Inverted Residuals and Linear Bottlenecks,'' in \emph{IEEE Conference on Computer Vision and Pattern Recognition (CVPR)}, pp.~4510--4520, 2018.
\bibitem{howard2019searching} A.~Howard et al., ``Searching for MobileNetV3,'' in \emph{IEEE/CVF International Conference on Computer Vision (ICCV)}, pp.~1314--1324, 2019.
\end{thebibliography}

\end{document}
